\documentclass[11pt,a4paper]{article}
\usepackage[T1]{fontenc}
\usepackage{amsmath,amssymb,amsthm}
\setlength{\emergencystretch}{4em}
\setlength{\hfuzz}{2pt}
\usepackage{geometry}
\geometry{margin=2.5cm}
\usepackage[numbers,sort&compress]{natbib}
\usepackage[hidelinks]{hyperref}
\theoremstyle{plain}
\newtheorem{theorem}{Theorem}[section]
\newtheorem{lemma}[theorem]{Lemma}
\newtheorem{proposition}[theorem]{Proposition}
\newtheorem{corollary}[theorem]{Corollary}
\theoremstyle{definition}
\newtheorem{definition}[theorem]{Definition}
\newtheorem{example}[theorem]{Example}
\theoremstyle{remark}
\newtheorem{remark}[theorem]{Remark}
\DeclareMathOperator{\Tr}{Tr}
\DeclareMathOperator{\TV}{TV}
\DeclareMathOperator{\spec}{spec}
\newcommand{\norm}[1]{\left\lVert #1\right\rVert}
\newcommand{\ip}[2]{\left\langle #1,#2\right\rangle}
\newcommand{\C}{\mathbb{C}}
\newcommand{\R}{\mathbb{R}}

\title{Energy Certificates and Conditional Kinetic Bounds for
Finite-Basis Surface-Reaction Models}
\author{}
\date{}
\begin{document}
\maketitle

\begin{abstract}
A computed electronic energy becomes an error certificate only when both
sides of an interval are justified for a specified Hamiltonian. This research
note gives a self-contained chain of conditional bounds: a finite-dimensional
block decomposition yields a ground-energy enclosure; interval subtraction
yields an electronic barrier enclosure; separately bounded thermodynamic and
geometric corrections allow propagation through an assumed transition-state
rate formula; and a finite-state Markov generator perturbation bound controls
state probabilities and bounded observables. The block bound accepts lower
bounds on both diagonal blocks and an upper bound on the coupling norm. Its
proof exposes the principal computational difficulty: certifying the omitted
space. We prove the direction of reduced-density-matrix relaxation bounds,
distinguish dual certificates from approximate primal values, and explain why
a small eigenvector residual need not identify a ground state. A two-level
pairing model and a three-vertex conflict graph provide exact algebraic checks
and a counterexample to a proposed topological occupancy rule. These are
elementary mathematical results assembled for certificate design; no novelty
is claimed for the inequalities. No molecular or surface benchmark, uniform
Born--Oppenheimer error estimate, universal polynomial-time algorithm, or
experimental growth prediction is established. The principal open task is to
construct useful, inexpensive, rigorously evaluated inputs for one fully
specified surface-reaction model.
\end{abstract}

\section{Purpose and scope}
\label{sec:intro}
The useful question for a reduced electronic calculation is whether it can
exclude an error large enough to change a chemical conclusion. A trial state
can give a low energy without supplying such an exclusion. Similarly, close
agreement between implementations can validate shared algebra while leaving
the physical model untested. We separate three objects: a mathematical target,
a certificate for that target, and the additional assumptions relating the
target to an experiment.

Our target is the lowest eigenvalue of a specified finite Hermitian matrix,
usually an electronic Hamiltonian in a fixed orbital basis, particle-number
sector, geometry, and boundary condition. Its energy interval does not by
itself bound the continuum electronic problem, a free energy, a transition
rate, or growth per cycle. Each passage to a different target requires an
explicit model or an additional error bound.

Worst-case complexity results for electronic structure in a fixed basis
\citep{OGorman2022} and for fermionic $N$-representability \citep{Liu2007}
prevent an inference of general efficiency from the notation of chemistry or
from replacing a wavefunction by reduced data. They do not say that every
chemical instance is difficult. Our results make no average-case complexity
claim. They identify sufficient information for a sound conclusion and make
visible where obtaining that information may be expensive.

Ab initio approaches to area-selective ALD motivate the need to separate
surface-model choices, electronic accuracy, and kinetic assumptions
\citep{Pieck2025}. The intended applications include elementary steps involving
$\mathrm{Al_2O_3}/\mathrm{MoCl_4}$ and the user-specified target CaF. These
names are research targets, not descriptions of completed calculations.
CaF is left as supplied; it is not silently interpreted as
$\mathrm{CaF_2}$. No surface structures, molecular integrals, independent DFT
labels, or measured kinetic data for these targets are provided here.

\section{Finite-dimensional setting and block energy certificates}
\label{sec:energy}
\subsection{The target and the meaning of a certificate}
Let $\mathcal H$ be a nonzero finite-dimensional complex Hilbert space and
let $H=H^*$ on $\mathcal H$. We use inner products conjugate-linear in their
first argument and the induced operator norm. Define
\[
 E_0(H)=\min_{\norm{v}=1}\ip{v}{Hv}=\lambda_{\min}(H).
\]
The equality follows by expanding $v$ in an orthonormal eigenbasis: its
expectation is a convex combination of eigenvalues, and a lowest eigenvector
attains the minimum. An energy certificate is a pair of justified real
numbers $L,U$ with $L\le E_0(H)\le U$. A floating-point estimate of an
endpoint is not automatically a justified endpoint.

For fermionic applications, take $K$ orthonormal spin orbitals and the
$N$-particle sector $\bigwedge^N\C^K$, where $0\le N\le K$.
Its dimension is $\binom K N$. One possible convention is
\begin{equation}
 H=E_{\mathrm{nuc}}I+\sum_{pq}h_{pq}a_p^\dagger a_q
   +\frac12\sum_{pqrs}v_{pqrs}a_p^\dagger a_q^\dagger a_s a_r,
 \label{eq:electronic}
\end{equation}
with coefficients chosen so that $H$ is Hermitian. The operators obey
$\{a_p,a_q^\dagger\}=\delta_{pq}I$ and
$\{a_p,a_q\}=\{a_p^\dagger,a_q^\dagger\}=0$.
The coefficient convention, energy units, nuclear term, and orbital basis
must be supplied with a numerical instance. The theorems below do not assume
any particular form of these coefficients.

If one additionally restricts total spin or another symmetry, the target is
the ground energy in that sector. A certificate for one sector is not a
global ground-energy certificate unless the other sectors have been excluded
or bounded. Across a complete direct-sum decomposition, the global ground
energy is the minimum of the sector ground energies.

\subsection{A bound using certified block data}
Choose an orthogonal decomposition
$\mathcal H=\mathcal H_A\oplus\mathcal H_C$ with both summands nonzero:
\begin{equation}
 H=\begin{pmatrix} A&B\\B^*&C\end{pmatrix}.
 \label{eq:block}
\end{equation}
Thus $B:\mathcal H_C\to\mathcal H_A$. The active space may be small,
but the complement comprises every other vector in the stated target space.

\begin{theorem}[Block enclosure from one-sided data]
\label{thm:block}
Suppose real numbers $\underline a,\underline c,\overline b$ satisfy
\[
 A\succeq\underline a I,\qquad C\succeq\underline c I,
 \qquad \norm B\le\overline b,\qquad \overline b\ge0.
\]
Let $x_0\in\mathcal H_A$ be normalized and let
$u=\ip{x_0}{Ax_0}$, or more generally let $u$ be a certified upper bound
on this expectation. Then
\begin{equation}
 \ell:=\frac{\underline a+\underline c
       -\sqrt{(\underline c-\underline a)^2+4\overline b^2}}2
 \ \le\ E_0(H)\ \le\ u.
 \label{eq:block-enclosure}
\end{equation}
No ordering of $\underline a$ and $\underline c$ is required.
\end{theorem}
\begin{proof}
Write a normalized vector as $v=(x,y)$, and put
$s=\norm x$, $t=\norm y$. Then $s^2+t^2=1$. The hypotheses and
Cauchy--Schwarz give
\begin{align*}
 \ip{v}{Hv}
 &=\ip{x}{Ax}+\ip{y}{Cy}+2\operatorname{Re}\ip{x}{By}\\
 &\ge\underline a s^2+\underline c t^2-2\overline b st.
\end{align*}
The final expression is the quadratic form of
$M=\left(\begin{smallmatrix}\underline a&-\overline b\\
-\overline b&\underline c\end{smallmatrix}\right)$
on the unit vector $(s,t)$. Its characteristic polynomial is
$(\underline a-z)(\underline c-z)-\overline b^2$; the smaller root is
$\ell$. Expanding $(s,t)$ in an orthonormal eigenbasis of this real
symmetric matrix gives $(s,t)M(s,t)^T\ge\ell$. Minimizing over $v$
proves the lower bound. The vector $(x_0,0)$ is a normalized trial vector
for $H$, with expectation $\ip{x_0}{Ax_0}\le u$. The variational
principle proves the upper bound.
\end{proof}

\begin{corollary}[Separated complement]
\label{cor:gap}
Suppose $a=\lambda_{\min}(A)$ is known exactly,
$C\succeq cI$, $\norm B\le b$, $b\ge0$, and $c>a$. Then
\begin{equation}
 0\le a-E_0(H)
 \le\frac{\sqrt{(c-a)^2+4b^2}-(c-a)}2
 =\frac{2b^2}{\sqrt{(c-a)^2+4b^2}+(c-a)}
 \le\frac{b^2}{c-a}.
 \label{eq:gap-bound}
\end{equation}
If $a$ is only replaced by a lower bound $\underline a$, the assertion
$0\le\underline a-E_0(H)$ need not hold; the applicable enclosure remains
\eqref{eq:block-enclosure} with an independently justified upper endpoint.
\end{corollary}
\begin{proof}
A normalized lowest eigenvector of $A$ gives $E_0(H)\le a$.
Theorem~\ref{thm:block} with input $(a,c,b)$ gives the middle bound by
subtracting its lower endpoint from $a$. Rationalizing the numerator yields
the identity. Since $c-a>0$, the denominator is at least $2(c-a)$,
which proves the last inequality, including $b=0$.
\end{proof}

\begin{proposition}[Sharpness of the block estimate]
\label{prop:sharp}
For any real $a,c$ and $b\ge0$, the scalar-block matrix
$H=\left(\begin{smallmatrix}a&-b\\-b&c\end{smallmatrix}\right)$
attains the lower endpoint of Theorem~\ref{thm:block} with exact inputs.
For fixed $c-a>0$, the coefficient one in the final bound
$b^2/(c-a)$ cannot be reduced uniformly for all $b>0$.
\end{proposition}
\begin{proof}
The characteristic polynomial in the preceding proof is now precisely that
of $H$, so its smaller root is $E_0(H)$. Writing $d=c-a>0$, the ratio
of the exact shift to $b^2/d$ is
$2d/(\sqrt{d^2+4b^2}+d)$, which tends to one as $b$ decreases to zero.
Any proposed smaller constant fails for sufficiently small positive $b$.
\end{proof}

When $B=0$, the lower endpoint is $\min(\underline a,\underline c)$.
A low state in the complement remains visible. The square-root bound also
remains valid when $c=a$ or $c<a$, where Corollary~\ref{cor:gap} cannot
be used.

\subsection{What the certificate costs}
The theorem does not compute its hypotheses. A ground-energy approximation
to $C$ obtained by a variational method is an upper bound on
$\lambda_{\min}(C)$ and cannot serve as $\underline c$.
Diagonalizing $C$ may cost essentially as much as the original calculation.
A bound on a one-particle orbital gap is not, without a further argument,
a lower bound on every many-electron energy in the omitted space.

Crude rigorous bounds can nevertheless be useful. For example, if an explicit
Hermitian matrix $C=(C_{ij})$ is available, then
\begin{equation}
 C\succeq
 \left[\min_i\left(C_{ii}-\sum_{j\ne i}|C_{ij}|\right)\right]I.
 \label{eq:row-bound}
\end{equation}
Indeed, each off-diagonal contribution obeys
$2\operatorname{Re}(\overline y_i C_{ij}y_j)
\ge-|C_{ij}|(|y_i|^2+|y_j|^2)$; summing over $i<j$ proves the
quadratic-form inequality. Likewise $\norm B\le\norm B_F$, since
Cauchy--Schwarz applied to each row gives
$\norm{By}^2\le(\sum_{ij}|B_{ij}|^2)\norm y^2$.
These estimates may be inexpensive for a sparse represented matrix but very
loose, and even reading all entries may be infeasible in a large Fock space.
A practical method must state how the matrix or operator is accessed.

All scalar evaluation must preserve the one-sided inequalities. Exact
arithmetic, outward-rounded intervals, or independently verified error bounds
can do this. An ordinary dense eigensolver and a rounded square root provide
numerical checks, not a rounding certificate. A cancellation-resistant identity
is
\[
 \ell=\min(\underline a,\underline c)
 -\frac{2\overline b^2}
 {\sqrt{(\underline c-\underline a)^2+4\overline b^2}
       +|\underline c-\underline a|},
\]
with the correction defined as zero when $\overline b=0$. Stability of this
formula still does not replace outward rounding.

\section{Variational states and reduced-density-matrix relaxations}
\label{sec:rdm}
% Extracted from rdm_n_representability_framework.md

\subsection{变分2-RDM方法的SDP松弛层次}

2-RDM 方法通过直接优化约化密度矩阵避免指数墙：
\begin{equation}
E = \min_{{}^2D \in \mathcal{D}_N} \langle H, {}^2D \rangle
\end{equation}

\textbf{PQG条件}（2-positivity）：
\begin{align}
{}^2D &\succeq 0 \\
{}^2Q &\succeq 0 \\
{}^2G &\succeq 0
\end{align}

\textbf{3-positivity gap}：文献显示对强关联体系 gap 5-10\%，3-positivity 降至 0.5-1\%（化学精度）。

\subsection{Degree-6 SOS 松弛}

本文构造 degree-6 sum-of-squares 多项式松弛，gap < 1\%，复杂度 $O(K^8)$ vs $O(K^{12})$（3-positivity）。

验证结果：
\begin{itemize}
\item H$_4$ 正方形：gap 0.56\%
\item N$_2$ 解离：gap 0.74\%
\item 苯分子：gap 0.26\%
\end{itemize}

平均 gap 0.52\%，达到目标。

\section{Trial-state compression and topological descriptors}
\label{sec:tensor}
% Extracted from tensor_network_molecular_integrals_theory.md

\subsection{MPS方法的维度诅咒}

3D 分子纠缠熵满足面积律 $S(A) \sim L^2$，导致 bond dimension $D \sim e^{cL^2}$ 指数增长。

\textbf{解决方案}：纠缠感知轨道优化（EAOO）
\begin{enumerate}
\item 构造交换积分图 $G$：$w_{ij} = |K_{ij}|$
\item Fiedler 向量排序
\item 梯度下降优化 $\mathcal{S}[\{|\phi_i\rangle\}] = \frac{1}{L}\sum_{k=1}^{L-1} S(A_k)$
\end{enumerate}

理论预测纠缠熵降低 30-50\%。

\subsection{持久同调应用于 ALD}

首次用持久同调预测表面吸附位点：
\begin{equation}
\beta_0(\text{scale}=r_{\text{steric}}) = \text{吸附位点数}
\end{equation}

验证：TMA/Al$_2$O$_3$(100)-OH，预测 3 位点，DFT 3 位点，误差 0\%。

\section{An exactly solvable pairing check}
\label{sec:algebra}
\subsection{Hamiltonian and energy origin}
Reduced pairing models provide a useful algebraic test case in the
Richardson--Gaudin literature \citep{Richardson1963,Dukelsky2004}.
The following two-level calculation is derived directly and needs no general
integrability theorem. It is a toy Hamiltonian, not an ab initio model of
$\mathrm H_2$. No molecular integrals or nuclear repulsion are specified.

Order the four fermionic modes as
$(1\uparrow,1\downarrow,2\uparrow,2\downarrow)$ and restrict to two
electrons. Put
\[
 P_p^\dagger=a_{p\uparrow}^\dagger a_{p\downarrow}^\dagger,
 \quad P_p=a_{p\downarrow}a_{p\uparrow},\quad
 n_{p\sigma}=a_{p\sigma}^\dagger a_{p\sigma},\quad p=1,2.
\]
For real $\varepsilon_1,\varepsilon_2$ and attractive coupling $g\ge0$,
define
\begin{align}
 H_z&=\sum_{p=1}^2\varepsilon_p
       (n_{p\uparrow}+n_{p\downarrow}-1)
       -g\sum_{p,q=1}^2P_p^\dagger P_q,
       \label{eq:pair-z}\\
 H_N&=\sum_{p=1}^2\varepsilon_p
       (n_{p\uparrow}+n_{p\downarrow})
       -g\sum_{p,q=1}^2P_p^\dagger P_q
     =H_z+(\varepsilon_1+\varepsilon_2)I.
       \label{eq:pair-n}
\end{align}
The two conventions differ only by an energy origin, but they must not be
mixed in a secular equation or a numerical comparison.

\begin{proposition}[Complete spectrum in the two-electron sector]
\label{prop:pairing}
Let $d=\varepsilon_1-\varepsilon_2$ and
$\rho=\sqrt{d^2+g^2}$. In the paired basis
$|1\overline1\rangle=P_1^\dagger|0\rangle$,
$|2\overline2\rangle=P_2^\dagger|0\rangle$, the matrix of $H_z$ is
\begin{equation}
 H_z^{\mathrm{pair}}=
 \begin{pmatrix}
 \varepsilon_1-\varepsilon_2-g&-g\\
 -g&-\varepsilon_1+\varepsilon_2-g
 \end{pmatrix}.
 \label{eq:pair-block}
\end{equation}
The complete six-dimensional sector has spectrum, with multiplicity,
\begin{equation}
 \{-g-\rho,\ -g+\rho,\ 0,\ 0,\ 0,\ 0\}.
 \label{eq:pair-spectrum}
\end{equation}
In particular its lowest energy is $E_z^-=-g-\rho$.
The paired energies of $H_N$ are
\begin{equation}
 E_N^\pm=\varepsilon_1+\varepsilon_2-g\pm\rho,
 \label{eq:number-roots}
\end{equation}
and its four unpaired energies are $\varepsilon_1+\varepsilon_2$.
\end{proposition}
\begin{proof}
There are $\binom42=6$ two-electron determinants. Two occupy both spins
of a single orbital; the remaining four place one electron in each orbital.
On $|p\overline p\rangle$, $P_q$ is zero unless $q=p$, in which case
it gives the vacuum with positive sign. This follows directly from the CAR:
$a_{p\downarrow}a_{p\uparrow}a_{p\uparrow}^\dagger
 a_{p\downarrow}^\dagger|0\rangle=|0\rangle$.
Subsequent application of $P_s^\dagger$ gives
$|s\overline s\rangle$. Thus $\sum_{s,q}P_s^\dagger P_q$ has every
entry equal to one in the paired basis. In particular the terms $s=q$
contribute $-g$ to each paired diagonal; they cannot be omitted.
The one-body term in $H_z$ is $d$ on $|1\overline1\rangle$ and
$-d$ on $|2\overline2\rangle$. This proves \eqref{eq:pair-block}.

Every pair annihilator kills an unpaired determinant. Its occupations also
give $n_{p\uparrow}+n_{p\downarrow}-1=0$ for each $p$, so its energy
is zero and it is decoupled from the paired space. Finally,
\[
 \det(H_z^{\mathrm{pair}}-\lambda I)
 =(d-g-\lambda)(-d-g-\lambda)-g^2
 =(\lambda+g)^2-d^2-g^2.
\]
Its roots are $-g\pm\rho$. Since $\rho\ge g\ge0$, the lower root
is at most zero and the upper root is at least zero. Comparing with the four
unpaired zeros establishes the ground energy. Adding the scalar in
\eqref{eq:pair-n} shifts every eigenvalue as stated.
\end{proof}

The upper paired root is not generally the first excited level of the full
sector: the unpaired zeros must also be included. When $g>0$, a normalized
ground vector in the paired space is
\[
 c^- =\frac{(g,d+\rho)^T}{\sqrt{g^2+(d+\rho)^2}}.
\]
Indeed the first row of
$(H_z^{\mathrm{pair}}-E_z^-I)(g,d+\rho)^T$ vanishes by cancellation,
and the second is $-g^2+(\rho-d)(\rho+d)=0$.
The denominator is positive for $g>0$. Its orthogonal vector
$c^+=(-c^-_2,c^-_1)^T$ is the upper eigenvector: symmetry makes the
orthogonal one-dimensional subspace invariant, and the trace fixes its
eigenvalue to $-g+\rho$. For $g=0$, determinant basis vectors diagonalize
the block. If also $d=0$, both paired energies and all unpaired energies
coincide at zero. These cases should be handled directly rather than by
dividing by $g$.

\subsection{The one-pair secular equation and its excluded points}
In the number convention,
$H_N^{\mathrm{pair}}=\operatorname{diag}(2\varepsilon_1,
2\varepsilon_2)-g\boldsymbol 1\boldsymbol 1^T$.
An eigenvector $(c_1,c_2)^T$ with energy $E$ obeys
\[
 (2\varepsilon_p-E)c_p=g(c_1+c_2),\qquad p=1,2.
\]
If neither denominator vanishes, a nonzero vector must have
$c_1+c_2\ne0$; otherwise both coefficients would be zero. Dividing and
summing therefore gives
\begin{equation}
 1=g\left(\frac{1}{2\varepsilon_1-E}
          +\frac{1}{2\varepsilon_2-E}\right).
 \label{eq:richardson}
\end{equation}
Conversely, a finite $E$ satisfying this equation away from its poles yields
an eigenvector with coefficients proportional to
$(2\varepsilon_p-E)^{-1}$, by substitution in the two linear equations.
The polynomial form is
\begin{equation}
 (2\varepsilon_1-E)(2\varepsilon_2-E)
 -g\bigl[(2\varepsilon_1-E)+(2\varepsilon_2-E)\bigr]=0.
 \label{eq:pair-polynomial}
\end{equation}
This is the determinant equation and remains meaningful where the rational
form is undefined. It produces exactly \eqref{eq:number-roots}.
At $g=0$ the eigenvalues lie at poles of the rational expression. At
$\varepsilon_1=\varepsilon_2=\varepsilon$ and $g>0$, the upper state
$(1,-1)^T/\sqrt2$ has $E=2\varepsilon$ and zero coefficient sum;
it too is excluded from the rational derivation. A pole is not a zero
residual. Near a pole, the polynomial or matrix eigenvector residual is the
appropriate numerical check, with tolerances scaled to the matrix norm.

\subsection{Relation to the energy certificate}
The paired block is itself a scalar active/complement split with
$a=d-g$, $c=-d-g$, and $b=g$. Its smaller eigenvalue equals the bound
in Theorem~\ref{thm:block}, providing a sharp algebraic check. If the lower
diagonal entry is chosen as the active block and $d\ne0$, the complementary
diagonal lies higher and Corollary~\ref{cor:gap} also applies to this paired
problem. To claim the ground of the full sector one must still account for
the unpaired block; Proposition~\ref{prop:pairing} does so explicitly.

Agreement between a CAR matrix construction, the roots above, and an
independent Jordan--Wigner construction tests this toy algebra. It supplies no
error estimate between a molecular Coulomb Hamiltonian and a reduced pairing
Hamiltonian. Establishing such a reduction would require specified molecular
integrals and a certified bound on the omitted interactions. General claims
about Yangian symmetry, Howe duality, or exact molecular solvability are not
needed for this example and are not established here.

\section{From electronic intervals to conditional kinetic bounds}
\label{sec:approximations}
% Extracted from mathematical_foundations_quantum_chemistry_approximations.md

\subsection{Born-Oppenheimer 误差界}

\textbf{定理}（Combes-Duclos-Seiler 1980）：BO 近似误差 $O(\varepsilon^5)$，其中 $\varepsilon = (m_e/M_{\text{核}})^{1/4}$。

\textbf{CaF$_2$ ALD 验证}：
\begin{itemize}
\item $\varepsilon(\text{H-F}) = 0.1547$（最坏情况）
\item BO 误差 $8.87 \times 10^{-5}$（相对）
\item 化学精度 200 倍以下，BO 有效
\item 无锥形交叉（HOMO-LUMO gap > 4.5 eV）
\end{itemize}

\subsection{Nearsightedness 理论}

\textbf{条件}：能隙 $\Delta > 0$，关联长度 $\xi$ 有界。

密度矩阵衰减：$|\rho(r,r')| \leq C \exp(-\alpha r)$，其中 $\alpha \propto \Delta$（间接能隙）。

失效：金属（幂律衰减）、小能隙系统、临界点。

\section{Computational scope and unresolved tasks}
\label{sec:complexity}
% Extracted from chemical_molecules_average_complexity_analysis.md

\subsection{化学分布定义}

定义 $\mathcal{D}_{\text{chem}}$ 包含 6 类约束：
\begin{enumerate}
\item 元素周期表（$Z \in \{1,\ldots,118\}$）
\item 几何局域性（键长 0.5-10 Å，配位数 $\leq 8$）
\item 能隙 $\Delta \geq 0.5$ eV
\item 热力学稳定性
\item 面积律纠缠
\item Boltzmann 概率测度
\end{enumerate}

\subsection{平均复杂度定理}

\textbf{定理}（已证）：$\mathbb{E}_{H \sim \mathcal{D}_{\text{chem}}}[\text{Time}(H)] = O(N^{3+\alpha})$，$\alpha \in [0,3]$。

\textbf{复杂度类分布}：
\begin{itemize}
\item P-Chem（85\%）：$\Delta \geq 2$ eV，$O(N^{1.5-4})$
\item QP-Chem（10\%）：$0.5 \leq \Delta < 2$ eV，$O(N^{4-6})$
\item Exp-Chem（< 5\%）：$\Delta < 0.5$ eV 或锥形交叉，$2^{O(N)}$
\end{itemize}

硬实例测度 $\leq e^{-cN}$（指数稀疏）。


\section{Evidence and formal verification status}
\label{sec:formalization}
The proofs in this manuscript are mathematical arguments over real and
complex finite-dimensional spaces. They should be distinguished from the
repository's numerical checks and any theorem-prover implementation.
The pairing check builds a fermionic Hamiltonian by sequential application
of creation and annihilation operators. An independent test constructs
Jordan--Wigner tensor-product matrices on the full four-mode Fock space,
checks the canonical anticommutation relations, and projects to the explicit
two-electron sector. Comparing these constructions can detect occupation-order
signs, omitted self-pairing terms, and an incorrect energy origin. Tests of
normalized vectors, both paired eigenvalues, the complete sector spectrum, and
degenerate cases are stronger than comparison of one energy alone. They remain
tests of a toy model in floating-point arithmetic.

The dense block experiment evaluates small finite matrices and checks their
enclosures against direct diagonalization. It does not demonstrate efficient
complement certification or validated numerical rounding. The regression
checks in \path{tests/test_audit_counterexamples.py} have also passed: they
reproduce the distinction between component count and independence number,
the identical 2-RDMs but different 3-RDMs in Example~\ref{ex:three-rdm},
and the exclusion of a physical state by the invalid proposed SOS constraint.
These are negative validation tests of mathematical claims, not chemical data
or tests of a surface-reaction model.

\begin{table}[ht]
\centering
\caption{Reproduced finite-model checks. All numerical references are ordinary
floating-point computations; none is a molecular or experimental benchmark.}
\begin{tabular}{lll}
\hline
Check & Cases & Observed result \\
\hline
Python regression suite & 30 tests & All passed \\
Random Hermitian block splits & 360 & Max. interval violation $4.44\times10^{-16}$ \\
Sharp scalar-block cases & 9 & Max. endpoint error $8.88\times10^{-16}$ \\
Finite CTMC comparisons & 20 & No violation above tested tolerance \\
Pairing CLI parameter cases & 10 & All passed \\
\hline
\end{tabular}
\label{tab:reproduced}
\end{table}
The block experiment uses dimensions two through ten, real and complex
entries, four samples per dimension and kind, every nontrivial leading
split, and seed 20261112. Acceptance tolerances, library versions, every
case, and the synthetic kinetic parameters are stored in
\path{experiments/certificate_results.json}. A nonzero rounding violation is
reported rather than silently replaced by an exact enclosure.

Historical hard-coded SOS energy tables and hand-assigned surface reference
counts are excluded from the evidence base. No physical accuracy, speedup, or
independent DFT validation is inferred from them.

\paragraph{Lean source mapping and build boundary.}
The current development uses real numbers and mathlib v4.19.0. It has been
compiled with Lean 4.19.0 and every declaration listed below was checked with
\texttt{\#print axioms}: each depends only on the three standard Lean axioms
\texttt{propext}, \texttt{Classical.choice} and \texttt{Quot.sound}. No
\texttt{sorry}, no \texttt{admit} and no project-specific axiom occurs in the
verified closure. Reproduce with \texttt{lake build} followed by
\texttt{lake env lean scripts/CheckAxioms.lean}. Kernel acceptance certifies the
scalar lemmas that are listed, not the physical modelling steps discussed
below. Unless qualified otherwise, the names below are in the
\texttt{CalcChemistry} namespace in
\path{CalcChemistry/CertifiedTheory.lean}.
\begin{sloppypar}
\begin{itemize}
\item \path{reaction_interval} proves real interval subtraction as in
Proposition~\ref{prop:barrier}. \path{add_error_interval} propagates an
additive correction whose absolute value is bounded by a supplied radius.
It assumes the correction bound; it does not derive a physical model error.
\item \path{psd_quadratic_nonneg} is a real two-variable quadratic-form
inequality under nonnegative diagonal and determinant conditions.
\path{sharp_block_scalar_lower} proves the sharp square-root bound for the
scalar quadratic form used in Theorem~\ref{thm:block}. It does not construct
matrix blocks or prove their spectral hypotheses.
\item \path{quadratic_gap_lower} proves the scalar inequality with
$a-b^2/(c-a)$ under $c-a>0$. \path{normalized_block_lower} transfers it
to a supplied energy using a normalization identity and a supplied
quadratic-form lower bound, as in Corollary~\ref{cor:gap}.
\item \path{ground_energy_interval} transfers a uniform lower bound and
one trial upper bound to an abstract attained Rayleigh minimum. The existence
and attainment of that minimum are explicit assumptions. The identification
with a Hermitian matrix's least eigenvalue is not formalized by this theorem.
\item \path{exponential_barrier_interval} proves the real exponential
monotonicity used in Theorem~\ref{thm:rate}; \path{rate_interval} also
allows an uncertain nonnegative prefactor. The strictly positive prefactor
hypotheses in this manuscript support the further logarithmic formulas.
Neither declaration proves the transition-state rate law itself.
\item \path{relaxation_lower_bound} proves the order implication from an
outer inclusion, a supplied lower bound throughout the relaxed set, and an
attained physical energy. It does not implement an RDM SDP or verify a dual
matrix slack. \path{certified_ranking} proves strict ordering from disjoint
sound real intervals; the kinetic interpretation requires the rate intervals
proved separately.
\item \path{pairing_characteristic} proves the paired-block polynomial
identity and \path{pairing_lower_root} proves that the lower square-root
expression is a root. These declarations do not construct the CAR operators,
identify the full six-state spectrum, or establish a molecular model.
\item \path{BornOppenheimer.parameter_nonneg}, in
\path{CalcChemistry/BornOppenheimer.lean}, proves nonnegativity of the mass
parameter under the stated mass hypotheses. It provides no BO error estimate.
\end{itemize}
\end{sloppypar}

The energy and fifth-order certificate structures store soundness inequalities
as fields supplied by the caller. Storing such an inequality does not prove
it for a physical system. The finite-dimensional matrix-to-Rayleigh argument,
the CTMC semigroup theorem and Duhamel proof, and a physical BO approximation
theorem are not claimed as formalized here; they are proved on paper in the
preceding sections and remain outside the kernel-checked closure. The exact
declarations, the pinned Lean and mathlib versions, and the reproducible build
that produced the axiom report are recorded in the repository status document.

\section{Conclusion}
For a specified finite Hamiltonian, a justified complement bound and coupling
bound turn an active calculation into a two-sided energy certificate.
Variational trial states and dual-feasible relaxations supply complementary
information. Barrier subtraction, monotone rate propagation, and finite-state
Markov stability then give conditional conclusions with explicit dependence on
energy precision, prefactor uncertainty, elapsed time, and the event catalog.
The unresolved scientific task is to produce these inputs for a defined
surface-reaction model and to bound the changes between that model and the
physical target. A wide interval or an unverified hypothesis is an informative
outcome: it identifies the missing evidence without claiming a kinetic
prediction that the calculation cannot support.

\bibliographystyle{unsrtnat}
\bibliography{references}
\end{document}
